\section{Finite rational presentations of Gaussian acquisitions}
\label{sec:v15-finite}
We now join finite resource certification to the actual marked physical
experiment. The intermediate model is allowed to have a nonfinite latent
state. Its reports, not merely its observable coordinates, must eventually
be finite. Throughout this section the parameter, action and mark alphabets
are finite, the horizon is finite, and the mark is drawn with the initial
state. An initial report, when present, is treated as one additional
acquisition. All channels below are parameter blind and task blind.

\subsection{A stateless report channel}
Let $(X,\mu)$ be a probability space, let the preparations be
$p_\theta\mu$ with $\|p_\theta\|_2\le D$, and let
$\gamma(dw\mid x)$ be the actual initial-mark channel. For each action
word $w\in A^k$, let $m_{k,w}\in L^2(\mu)$ be a real mean function.
The report at time $k$ is
\[
 Y_k=m_{k,a_{1:k}}(x)+\sigma_k\xi_k,\qquad \sigma_k>0,
\]
where the noises are independent standard Gaussians, independent of
$(x,W)$. Actions may be selected from all previously emitted reports.
The finite-dimensional models in the preceding section have this form;
so do the enriched approximations in Section~\ref{sec:v15-enriched}.
Put
\[
 B_k=\left(\sum_{w\in A^k}\|m_{k,w}\|_2^2\right)^{1/2}.
\]
Choose $R>0$ and a finite interval partition of $[-R,R]$ with maximum
cell length $h$. Adjoin the two tail cells. Let $Q$ return the cell of
a real report. Let $H$ reconstruct an interior cell uniformly and
reconstruct the tail cells at $-R$ and $R$, respectively. The endpoints
may be assigned arbitrarily to adjacent cells. These maps use no
persistent memory, seed, parameter or mark. Fresh randomness in $H$
is transient. Denote by $F^{R,h}$ the physical acquisition with the
report $Q(Y_k)$; its latent state and actual mark are unchanged.

\begin{lemma}[Uniform marked report quantization]\label{lem:v15-quantization}
The channel $Q$ simulates $F^{R,h}$ from $F$ exactly. The reverse
channel $H$ has marked feedback error at most
\begin{equation}\label{eq:v15-quantization}
 \beta(R,h)=1\wedge\sum_{k=1}^T
 \left\{\frac{h}{\sigma_k\sqrt{2\pi}}
       +\frac{2DB_k}{R}+\frac{4\sigma_k^2}{R^2}\right\}.
\end{equation}
The comparison includes the joint law of every independent exposed
selector and is uniform over every finite consumer width. In particular
$h\downarrow0$ and $R\uparrow\infty$ give two-sided stateless marked
comparison with error tending to zero.
\end{lemma}
\begin{proof}
For an absolutely continuous density $g$ and an interval $I$ of length
at most $h$, write $g_I=|I|^{-1}\int_Ig$. The fundamental theorem of
calculus and Fubini give
\[
 \int_I|g-g_I|\le |I|^{-1}\int_{I^2}|g(u)-g(v)|\,du\,dv
                  \le h\int_I|g'|.
\]
Consequently the variation error of the interior histogram is at most
$\frac h2\int|g'|$. For a Gaussian of variance $\sigma^2$ the latter
integral is $\sqrt{2/\pi}/\sigma$. Replacement of the tails costs at
most their total probability, even though the reconstruction is atomic.
For any mean $m$, this probability is at most
\[
 \1_{\{|m|>R/2\}}+\Pp\{|\sigma\xi|>R/2\}
 \le \frac{2|m|}{R}+\frac{4\sigma^2}{R^2}.
\]
Thus the conditional report error, given the original $x$ and the
selected action word, has the three terms in
\eqref{eq:v15-quantization} before integration.

Couple the exact and reconstructed systems with the same original
$(x,W)$ and independent selector. Before their first emitted-report
mismatch use the same controller randomness and maximally couple the
next reports. Matched-history weights of different words at time $k$
sum to at most one. Therefore the mean term is bounded by
\[
 \frac2R\int p_\theta(x)\max_{w\in A^k}|m_{k,w}(x)|\,d\mu(x)
 \le \frac{2D}{R}\left(\sum_w\|m_{k,w}\|_2^2\right)^{1/2}.
\]
This uses the original preparation, not an assumed bound for a
feedback-conditioned posterior. Sum the first-mismatch bounds.
On the matching event the complete emitted transcript, controls,
mark and selector agree. This proves the reverse comparison. Applying
$Q$ to each original report is the exact forward simulation, also
under feedback. Conditioning on the exposed selector throughout proves
the joint assertion. Composition never enlarges a retained register.
\end{proof}

The bound is deliberately an elementary certified bound; Gaussian tail
integrals or higher moments may improve it. No uniform bound on the
pointwise means is needed. The quantity $B_k$ can grow with the number
of action words. Uniformity in memory width is not uniformity in the
size of the finite instrument being constructed.

\subsection{Rationalization without a small-prefix denominator}
Let $G$ now be a finite-report acquisition of this form. For each
$\theta$, write
\[
 \pi_\theta(w)=\Pp_\theta(W=w),\qquad
 p_\theta(w,h;y\mid a)
       =\Pp_\theta(W=w,\,Y_{1:k-1}=h_y,\,Y_k=y
                         \mid a_{1:k-1}=h_a,a_k=a)
\]
for its joint marked prefix masses under the indicated open-loop
control word. Here $h=(h_y,h_a)$, and the conditioning on controls means
intervention, not conditioning on the probability of choosing a control.
Causality and the initial-mark convention imply
\[
 \sum_y p_\theta(w,h;y\mid a)=p_\theta(w,h)
\]
independently of the current action $a$. A parameterized latent
history-tree instrument may use these masses; its latent dependence
on $\theta$ is part of the experiment, not information given to a
simulator.

Suppose nonnegative rational arrays $t_\theta(w)$ and
$t_\theta(w,h;y\mid a)$ approximate the indicated joint masses.
Normalize $t_\theta(w)$ to obtain a rational initial-mark law. At each
node $(w,h,a)$ normalize its vector of children to obtain a rational
report kernel. If a normalizing sum is zero, use a fixed rational
probability vector, for example the uniform vector. This defines a
coherent finite rational marked causal experiment $G_{\mathbb Q}$.
Its prefix probabilities need not equal the originally rounded joint
masses; they are the probabilities generated by these normalized
kernels. That distinction is essential.

\begin{lemma}[Weighted normalization]\label{lem:v15-normalization}
If $u,v\in[0,\infty)^b$, $s=\sum u_j$, and $\bar u=u/s$ when $s>0$,
then, for $\bar v=v/\sum v_j$ with an arbitrary probability vector
chosen when its denominator vanishes,
\begin{equation}\label{eq:v15-normalization}
 s\|\bar u-\bar v\|_{\rm TV}\le\sum_j|u_j-v_j|.
\end{equation}
For $s=0$ the left side is defined to be zero.
\end{lemma}
\begin{proof}
If both sums are positive, the triangle inequality gives
\[
 \sum_j|u_j-s\bar v_j|
 \le\sum_j|u_j-v_j|+|\sum_jv_j-s|
 \le2\sum_j|u_j-v_j|.
\]
If $\sum v_j=0$, the left side of \eqref{eq:v15-normalization} is
at most $s=\sum|u_j-v_j|$. The remaining case is immediate.
\end{proof}

\begin{theorem}[Rational causal tree bridge]\label{thm:v15-rational}
The identity report channels compare $G$ and $G_{\mathbb Q}$ in both
directions, including the actual marks, with feedback error at most
\begin{align}\label{eq:v15-rational}
 \rho=1\wedge\max_\theta\bigg[&\sum_w|\pi_\theta(w)-t_\theta(w)|\\[-2mm]
 &+\sum_{k=1}^T\sum_{w,h,a,y}
       |p_\theta(w,h;y\mid a)-t_\theta(w,h;y\mid a)|\bigg].\nonumber
\end{align}
Here the finite sum includes all open-loop histories at the specified
time, not only the histories reached by a particular controller. The
bound is independent of simulator and consumer memory. It remains
valid with a visible selector and its joint reference law.

If every joint mass has computable rational enclosures of arbitrarily
small width, the construction is effective and produces $\rho\to0$.
No lower bound for any nonzero prefix probability, and no decision
whether a prefix has exactly zero probability, is required.
\end{theorem}
\begin{proof}
Couple the initial marks by their normalized laws. Lemma~\ref{lem:v15-normalization}
bounds the initial failure probability by the first sum in
\eqref{eq:v15-rational}. Use the same controller and selector on both
sides. Conditional on a matched marked history, maximally couple the
next reports. Under the original instrument, the probability of
reaching a particular marked history under any deterministic controller
is either zero or the corresponding open-loop prefix mass; for a
randomized controller it is at most that mass. The probability of
reaching it with the coupling still matched is smaller still. Thus
its first-mismatch contribution is bounded by the original prefix
mass times the conditional variation distance. Applying
\eqref{eq:v15-normalization} at that node bounds this product by the
sum of joint child-mass errors. Sum over all nodes and times. The
coupling inequality gives the asserted joint variation distance.
The estimate is symmetric as a comparison of the two emitted laws;
identity channels require no state in either direction.

For effectiveness, take a nonnegative rational lower endpoint of a
certified enclosure for every mass, clipping a negative lower endpoint
at zero. If there are $M$ mass entries for a parameter and every
interval has width at most $\eta/M$, the displayed bound is at most
$\eta$. Normalization uses only rational arithmetic. A zero rounded
sum invokes the declared fallback. The weighted bound, rather than a
bound on a ratio at an individual rare history, proves convergence.
Conditioning on an independent visible selector leaves every bound
unchanged. The mark is part of the finite tree and is never redrawn
independently of its reports.
\end{proof}

For the Gaussian acquisition its required masses are explicitly
\begin{equation}\label{eq:v15-mass-integral}
 \int_X p_\theta(x)\gamma(w\mid x)
       \prod_{j=1}^k
       \Pp\{m_{j,a_{1:j}}(x)+\sigma_j\xi_j\in I_{y_j}\}\,d\mu(x).
\end{equation}
This is an integration problem for a specified preparation and trial
family. The theorem does not infer its effective solvability from the
mere words ``standard Borel''. Rigorous enclosures for
\eqref{eq:v15-mass-integral} are part of the effective data specification
in Theorem~\ref{thm:v15-main}; explicit physical families are treated in
Section~\ref{sec:v15-physical-data}.
